\item Un motor funciona en un ciclo de Carnot de la siguiente manera. En el punto $A$ del ciclo, 2.34 moles de un gas ideal monoatómico tiene una presión de $1\,400\text{ kPa}$, un volumen de $10.0\text{ L}$ y una temperatura de $720\text{ K}$. El gas se expande isotrópicamente al punto $B$ y luego se expande adiabáticamente al punto $C$, donde su volumen es de $24.0\text{ L}$. Una compresión isotérmica lo lleva al punto $D$, donde su volumen es de $15.0\text{ L}$. Un proceso adiabático devuelve el gas al punto $A$. 
    
(a) Complete la siguiente tabla con presiones, volúmenes y temperaturas en cada uno de los cuatro puntos en el ciclo:

\begin{center}
\begin{tabular}{lccc}
\hline
\textbf{Punto en el ciclo} & \textbf{$P$(kPa)} & \textbf{$V$(L)} & \textbf{$T$(K)} \\
\hline
$A$ & & & \\
$B$ & & & \\
$C$ & & & \\
$D$ & & & \\
\hline
\end{tabular}
\end{center}

(b) Complete la siguiente tabla con la transferencia de energía por calor, y trabaje en el gas y el cambio en la energía interna del gas para cada uno de los cuatro procesos en el ciclo:

\begin{center}
\begin{tabular}{lccc}
\hline
\textbf{Punto en el ciclo} & \textbf{$Q$ (kJ)} & \textbf{$W$ (kJ)} & \textbf{$\Delta E_{\text{int}}$ (kJ)} \\
\hline
$A \rightarrow B$ & & & \\
$B \rightarrow C$ & & & \\
$C \rightarrow D$ & & & \\
$D \rightarrow A$ & & & \\
\hline
\end{tabular}
\end{center}

(c) Encuentre la eficiencia del motor a partir de los datos en la tabla en el inciso (b). 
(d) Encuentre la eficiencia del motor a partir de los datos de la tabla en el inciso (a).
